\section{21 条二次约束的完备性：维数计数}\label{app:dimcount}

令 $\nu_2:\mathbb{RP}^3\to\mathbb{RP}^9$ 为二次 Veronese 嵌入 $[q]\mapsto[v_2(q)]$，其中 $v_2(q)\in\mathbb{R}^{10}$ 按第~\ref{sec:lift}~节的归一化收集 $q=(q_0,q_1,q_2,q_3)$ 的十个二次单项式；并令
\begin{equation}
\begin{split}
&\mathcal{I}_2\big(\nu_2(\mathbb{RP}^3)\big)\\
&\quad=\Big\{M\in\operatorname{Sym}^2(\mathbb{R}^{10}):\\
&\qquad v_2(q)^\top Mv_2(q)\equiv0\quad\forall q\in\mathbb{R}^4\Big\}.
\end{split}
\label{eq:app-ideal}
\end{equation}
考虑线性求值映射
\begin{equation}
\begin{split}
&\Phi:\operatorname{Sym}^2(\mathbb{R}^{10})\longrightarrow\mathbb{R}[q_0,q_1,q_2,q_3]_4,\\
&\Phi(M)(q):=v_2(q)^\top Mv_2(q).
\end{split}
\label{eq:app-phi}
\end{equation}

\textbf{(i) 核。} 由 \eqref{eq:app-ideal}，$\ker\Phi=\mathcal{I}_2(\nu_2(\mathbb{RP}^3))$。

\textbf{(ii) 满射性。} 目标空间由满足 $a_0+a_1+a_2+a_3=4$ 的四次单项式 $q_0^{a_0}q_1^{a_1}q_2^{a_2}q_3^{a_3}$ 张成。每个这样的单项式可分解为两个二次单项式之积 $m_1(q)m_2(q)$：取整数 $0\le b_i\le a_i$ 且 $\sum_ib_i=2$（因 $\sum_ia_i=4$ 故可行），令 $m_1=\prod_iq_i^{b_i}$、$m_2=\prod_iq_i^{a_i-b_i}$。每个 $m_\ell$ 至多差一个常数就是 $v_2$ 的某个坐标；故存在指标 $r,s$ 与常数 $c_r,c_s$ 使
\begin{equation}
\begin{split}
m_1(q)m_2(q)&=c_rc_s\,(v_2)_r(q)(v_2)_s(q)\\
&=\Phi\!\left(\tfrac{c_rc_s}{2}(e_re_s^\top+e_se_r^\top)\right)(q).
\end{split}
\label{eq:app-monomial}
\end{equation}
于是像包含全部基单项式，$\Phi$ 为满射。

\textbf{(iii) 计数。} 源空间与目标空间的维数分别为
\begin{equation}
\begin{split}
\dim\operatorname{Sym}^2(\mathbb{R}^{10})&=\frac{10\cdot11}{2}=55,\\
\dim\mathbb{R}[q_0,q_1,q_2,q_3]_4&=\binom{4+4-1}{4}=\binom74=35 .
\end{split}
\label{eq:app-dims}
\end{equation}
由秩-零化度定理与满射性，
\begin{equation}
\dim\mathcal{I}_2\big(\nu_2(\mathbb{RP}^3)\big)=\dim\ker\Phi=55-35=20 .
\label{eq:app-20}
\end{equation}
这是关于映射 $\Phi$ 的陈述；它对于 21 条 $SO(3)$ 约束矩阵的完备性是必要而不充分的，因为尚未说明这些矩阵达到该计数。

\textbf{(iv) 21 条约束达到 20 维。} 设 $A_i\in\mathbb{S}^{10}$（$i=1,\dots,21$）为定理~\ref{thm:ideal} 的约束在 $x_R$ 坐标下的矩阵，$N_i:=(2U)^\top A_i(2U)$ 为其在 $v_2$ 坐标下的拉回。由于 $x_R=2Uv_2$ 且 $x_R(q)^\top A_ix_R(q)\equiv0$ 在 $\mathbb{S}^3$ 上成立，
\begin{equation}
v_2(q)^\top N_iv_2(q)=x_R(q)^\top A_ix_R(q)\equiv0\quad\forall q\in\mathbb{R}^4,
\label{eq:app-pullback}
\end{equation}
故 $N_i\in\ker\Phi$。对 $N_1,\dots,N_{21}$ 的 55 个系数向量作精确有理高斯消元，得
\begin{equation}
\operatorname{rank}\{N_1,\dots,N_{21}\}=20 .
\label{eq:app-rank}
\end{equation}
唯一的线性相关是迹恒等式
\begin{equation}
\begin{split}
&\sum_{a=1}^{3}\big(\|c_a\|^2-x_0^2\big)-\sum_{a=1}^{3}\big(\|r_a\|^2-x_0^2\big)\\
&\quad=\operatorname{tr}(R^\top R)-\operatorname{tr}(RR^\top)\equiv0 .
\end{split}
\label{eq:app-dependency}
\end{equation}
因此 $\operatorname{span}\{N_i\}$ 是 20 维核 \eqref{eq:app-20} 的 20 维子空间，故等于该核，21 条约束张成完整的二次消没理想。该计数对约束集敏感：仅六个非对角内积的秩为 $15$，这正是已发布 24 约束实现的约束亏缺并非记账细节的原因。

\section{Haar 测度、双覆盖与 Laplace 常数}\label{app:haar}

\subsection*{局部 Jacobian \texorpdfstring{$dR=8\,du$}{dR=8du}}

设 $q^{(k)}\in\mathbb{S}^3$，$E_k\in\mathbb{R}^{4\times3}$ 满足 $E_k^\top q^{(k)}=0$、$E_k^\top E_k=I_3$。取正规坐标
\begin{equation}
q(u)=\frac{q^{(k)}+E_ku}{\sqrt{1+\|u\|_2^2}},\qquad u\in\mathbb{R}^3 .
\label{eq:app-chart}
\end{equation}
在 $u=0$ 处 $\dot q(0)=E_k\dot u$，且体坐标系纯虚四元数为 $(q^{(k)})^{-1}\dot q(0)=[0,\dot u^\top]^\top$。在双覆盖 $\pi:\mathbb{S}^3\to SO(3)$（$R(q)v=qvq^{-1}$）下求导得
\begin{equation}
\mathrm{d}\pi_{q^{(k)}}\big(E_k\dot u\big)=R(q^{(k)})\,[2\dot u]_\times ,
\label{eq:app-dpi}
\end{equation}
故体角速度为 $\omega=2\dot u$。对一般标准正交切基 $E_k=E_R(q^{(k)})Q_k$（$Q_k\in O(3)$），有 $\partial\omega/\partial u|_{u=0}=2Q_k$。取与轴角向量欧氏测度对应的双不变体积元 $\mathrm{d}R=\mathrm{d}\omega_1\wedge\mathrm{d}\omega_2\wedge\mathrm{d}\omega_3$，则
\begin{equation}
\mathrm{d}R=\left|\det\frac{\partial\omega}{\partial u}\right|_{u=0}\mathrm{d}u
=|\det(2Q_k)|\,\mathrm{d}u=8\,\mathrm{d}u,
\label{eq:app-jac}
\end{equation}
且 $\operatorname{Vol}(SO(3))=8\pi^2$。
全局归一化一致：$\int_{SO(3)}f\,\mathrm{d}R=\frac12\int_{\mathbb{S}^3}f(R(q))\,8\,\mathrm{d}\sigma(q)=4\int_{\mathbb{S}^3}f\,\mathrm{d}\sigma$，取 $f=1$ 得 $\frac12\cdot8\cdot2\pi^2=8\pi^2$。

\subsection*{双覆盖与逐模态因子}

在非退化模态 $q^{(k)}$ 附近，正规坐标 \eqref{eq:app-chart} 中单个叶片上的 Laplace 积分为 $\int e^{-J/\tau}\mathrm{d}u=(2\pi\tau)^{3/2}(\det H_k)^{-1/2}e^{-J_k/\tau}(1+O(\tau))$。乘 Jacobian \eqref{eq:app-jac} 得每叶片 $8(2\pi\tau)^{3/2}(\det H_k)^{-1/2}e^{-J_k/\tau}$。对映两片 $\pm q^{(k)}$ 是 $\mathbb{S}^3$ 上 $J$ 与 $H$ 相同的两个不同点，双覆盖给出因子 $1/2$；两个因子精确抵消，故每个 $SO(3)$ 模态贡献
\begin{equation}
8(2\pi\tau)^{3/2}\frac{e^{-J_k/\tau}}{\sqrt{\det H_k}} .
\label{eq:app-permode}
\end{equation}
于是，在未归一化测度 $\mathrm{d}R$ 下，
\begin{equation}
\begin{split}
&\int_{SO(3)}e^{-J_{\mathrm{red}}/\tau}\mathrm{d}R\\
&\quad=8(2\pi\tau)^{3/2}\sum_{k=1}^K\frac{e^{-J_k/\tau}}{\sqrt{\det H_k}}\big(1+O(\tau)\big),
\end{split}
\label{eq:app-rotint}
\end{equation}
即 \eqref{eq:rotationint}。若取概率 Haar 测度 $\mathrm{d}\mu_{\mathrm{Haar}}=\mathrm{d}R/(8\pi^2)$，则同一积分再除以 $8\pi^2$。

\subsection*{\texorpdfstring{$Z_{\mathrm{Lap}}$}{Z\_Lap} 的绝对常数}

将 \eqref{eq:app-rotint} 与精确平移积分 \eqref{eq:gaussint} 结合，并取概率 Haar 测度，
\begin{equation}
\begin{split}
Z_{\mathrm{Lap}}(\theta)&=\frac{(\pi\tau)^{3/2}}{\sqrt{\det G}}\cdot\frac{(2\pi\tau)^{3/2}}{\pi^2}\sum_{k=1}^K\frac{e^{-J_k/\tau}}{\sqrt{\det H_k}}\\
&=\frac{2\sqrt2\,\pi\tau^3}{\sqrt{\det G}}\sum_{k=1}^K\frac{e^{-J_k/\tau}}{\sqrt{\det H_k}} .
\end{split}
\label{eq:app-zlap}
\end{equation}
若取未归一化测度 $\mathrm{d}R$，右端还要乘与 $\theta$ 无关的因子 $8\pi^2$。该因子在固定 $\tau$ 时为常数，在 $w_k$ 与 $\nabla_\theta L_{\tau\text{-KL}}$ 中相消；只有 \eqref{eq:app-zlap} 应作为绝对常数引用。

\subsection*{包络定理与 \texorpdfstring{$\log\det H_k$}{log det H\_k} 的完整导数}

设 $J_{\mathrm{red}}(q;\theta)=v_2(q)^\top M(\theta)v_2(q)$，$q^{(k)}(\theta)$ 为非退化驻点，$E_k^\top\nabla_qJ_{\mathrm{red}}=0$，$H_k=E_k^\top(\nabla^2J_{\mathrm{red}}-4J_kI_4)E_k\succ0$。由于 $\mathrm{d}q^{(k)}\in\operatorname{range}(E_k)$ 且 $\operatorname{range}(E_k)=(q^{(k)})^\perp$，
\begin{equation}
\nabla_\theta J_k=v_2(q^{(k)})^\top\big(\partial_\theta M\big)v_2(q^{(k)}) .
\label{eq:app-envelope}
\end{equation}
记 $\xi_k=-H_k^{-1}E_k^\top\partial_\theta\nabla_qJ_{\mathrm{red}}$，使 $\partial_\theta q^{(k)}=E_k\xi_k$，由链式法则得
\begin{equation}
\begin{split}
&\partial_\theta\log\det H_k\\
&\quad=\operatorname{tr}\!\Big[H_k^{-1}E_k^\top\Big(\partial_\theta\nabla^2J_{\mathrm{red}}\\
&\qquad+\mathrm{D}_q\big(\nabla^2J_{\mathrm{red}}\big)[E_k\xi_k]
-4\,\partial_\theta J_k\,I_4\Big)E_k\Big].
\end{split}
\label{eq:app-dlogdet}
\end{equation}
有两个精确简化。反对称标架旋转项在 $\operatorname{tr}(H_k^{-1}\dot H_k)$ 中由迹的循环性相消，因为 $\log\det H_k$ 在切基的正交变换下不变。法向弯曲项由 Euler 齐次性消失：$\nabla^2J_{\mathrm{red}}q^{(k)}=3\nabla_qJ_{\mathrm{red}}(q^{(k)})$，故 $E_k^\top\nabla^2J_{\mathrm{red}}q^{(k)}=3E_k^\top\nabla_qJ_{\mathrm{red}}=0$。最后，$J_{\mathrm{red}}$ 是四次型，故 $\mathrm{D}_q(\nabla^2J_{\mathrm{red}})[\delta q]$ 关于 $q$ 线性，\eqref{eq:app-dlogdet} 中的缩并每个模态仅需 $O(1)$。这就是简写 \eqref{eq:dlogdet} 与梯度 \eqref{eq:annealedgrad} 背后的完整导数。
